% Aim: Bulleted, numbered and definition lists; compact/customised lists; interrupting and resuming
\documentclass[12pt]{article}
\usepackage{enumitem}   % for customising lists

\begin{document}

\section*{Lists in \LaTeX{}}

\subsection*{1. Bulleted list (itemize)}
\begin{itemize}
  \item Algebra
  \item Calculus
  \item Real Analysis
\end{itemize}

\subsection*{2. Numbered list (enumerate)}
\begin{enumerate}
  \item Open the editor
  \item Write the code
  \item Compile the file
\end{enumerate}

\subsection*{3. Definition list (description)}
\begin{description}
  \item[Group] A set with an operation satisfying closure, associativity, identity and inverse.
  \item[Ring] A set with two operations, addition and multiplication.
  \item[Field] A commutative ring in which every non-zero element has an inverse.
\end{description}

\subsection*{4. Compact list (spacing removed)}
\begin{itemize}[noitemsep]
  \item Matrices
  \item Determinants
  \item Vectors
\end{itemize}

\subsection*{5. Customised spacing}
\begin{enumerate}[itemsep=10pt, topsep=5pt]
  \item First item with extra space
  \item Second item with extra space
  \item Third item with extra space
\end{enumerate}

\subsection*{6. Customised labels}
\begin{enumerate}[label=(\alph*)]
  \item Sets
  \item Relations
  \item Functions
\end{enumerate}

\begin{itemize}[label=$\star$]
  \item Star bullet one
  \item Star bullet two
\end{itemize}

\subsection*{7. Nested list}
\begin{enumerate}
  \item Mathematics
    \begin{itemize}
      \item Pure
      \item Applied
    \end{itemize}
  \item Computer Science
\end{enumerate}

\subsection*{8. Interrupting and resuming a list}
\begin{enumerate}
  \item First point
  \item Second point
\end{enumerate}
This is a normal paragraph in between. The list is interrupted here and will be resumed below.
\begin{enumerate}[resume]
  \item Third point (numbering continues)
  \item Fourth point
\end{enumerate}

\end{document}
